\documentclass[reqno]{amsart}
\usepackage[T1]{fontenc}
\usepackage{amsmath,amssymb,amsthm}
\usepackage{xr-hyper}
\usepackage[hidelinks]{hyperref}
\externaldocument[ds-]{20261125T004000--denominators-and-sha__workbook}
\externaldocument[mi-]{20261126T020430--the-minus-part-by-gauss-sums__workbook}
\externaldocument[da-]{20261126T030300--the-data__workbook}

% Chapter number; \ref{chap} gives it.
\newcounter{chap}
% Chapter numbers, generated from 20261125T012000--grh-workbook__guide.org
% by scripts/make-chapters.py.  Do not edit: change the order of the
% index and rerun the script (./preview-site.sh and CI do this).
%
% Usage in a sheet, after \newcounter{chap}:
%   % Chapter numbers, generated from 20261125T012000--grh-workbook__guide.org
% by scripts/make-chapters.py.  Do not edit: change the order of the
% index and rerun the script (./preview-site.sh and CI do this).
%
% Usage in a sheet, after \newcounter{chap}:
%   % Chapter numbers, generated from 20261125T012000--grh-workbook__guide.org
% by scripts/make-chapters.py.  Do not edit: change the order of the
% index and rerun the script (./preview-site.sh and CI do this).
%
% Usage in a sheet, after \newcounter{chap}:
%   \input{chapters}\chapnumber{<sheet ID>}
\makeatletter
\providecommand\chapters@set[2]{\@namedef{chapters@#1}{#2}}
\chapters@set{20261125T001000}{1}% zeros-to-ramp-sums
\chapters@set{20261125T001100}{2}% inert-primes-and-estimate-b
\chapters@set{20261125T005000}{3}% galois-fourier-and-averaging
\chapters@set{20261125T005500}{4}% columns-from-a-sphere
\chapters@set{20261125T000000}{5}% slope-method
\chapters@set{20261126T030100}{6}% a-twist-for-each-prime
\chapters@set{20261125T003000}{7}% the-chain-and-the-square-root
\chapters@set{20261125T010000}{8}% rank-lemma-and-zero-estimate
\chapters@set{20261125T010500}{9}% poincare-series-and-near-orthogonality
\chapters@set{20261124T223000}{10}% elliptic-curves-for-the-proof
\chapters@set{20261125T009500}{11}% fock-model-and-theta-intertwiners
\chapters@set{20261124T223100}{12}% complex-multiplication-and-inert-reduction
\chapters@set{20261125T008000}{13}% one-cm-determinant
\chapters@set{20261124T214500}{14}% supersingular-theta-divisibility
\chapters@set{20261124T214600}{15}% the-isotropic-plane
\chapters@set{20261125T013700}{16}% pairings-on-a-formal-group-of-height-two
\chapters@set{20261125T013500}{17}% hecke-nondegeneracy-on-the-tree
\chapters@set{20261126T030200}{18}% the-bound-is-attained
\chapters@set{20261125T009000}{19}% kronecker-limit-and-chowla-selberg
\chapters@set{20261125T002000}{20}% definite-theta-lifts
\chapters@set{20261125T015500}{21}% siegel-weil-in-the-smallest-case
\chapters@set{20261125T002100}{22}% definite-theta-lifts-and-norms
\chapters@set{20261125T016500}{23}% hecke-operators-and-theta-lifts
\chapters@set{20261125T016700}{24}% divisibility-along-the-chain
\chapters@set{20261125T004000}{25}% denominators-and-sha
\chapters@set{20261125T004100}{26}% tate-twists-and-sha
\chapters@set{20261126T020430}{27}% the-minus-part-by-gauss-sums
\chapters@set{20261125T014100}{28}% congruence-modules
\chapters@set{20261125T014500}{29}% eisenstein-congruences-and-selmer-classes
\chapters@set{20261125T014700}{30}% pseudocharacters-and-the-length-bound
\chapters@set{20261125T011000}{31}% galois-representations-black-box
\chapters@set{20261126T030300}{32}% the-data
\chapters@set{20261126T020210}{33}% bessel-poincare-series-and-their-gram-operators
\chapters@set{20261126T030400}{34}% the-seesaw-identity
\chapters@set{20261126T030500}{35}% divisibility-of-the-composite-lift
\chapters@set{20261126T030600}{36}% theorem-a-short-forms
\chapters@set{20261126T030700}{37}% the-determinant-of-a-chain
\chapters@set{20261126T030800}{38}% the-cost-of-one-lift
\chapters@set{20261126T030900}{39}% denominators-of-projections-to-lifts
\chapters@set{20261126T031000}{40}% twists-with-few-congruences
\chapters@set{20261126T031100}{41}% theorem-c-and-the-main-inequality
\providecommand\chapnumber[1]{%
  \@ifundefined{chapters@#1}%
    {\PackageError{chapters}{Sheet #1 is not listed in the index}%
      {List it in 20261125T012000--grh-workbook__guide.org\MessageBreak
       and run scripts/make-chapters.py.}}%
    {\edef\thechap{\@nameuse{chapters@#1}}}}
\makeatother
\chapnumber{<sheet ID>}
\makeatletter
\providecommand\chapters@set[2]{\@namedef{chapters@#1}{#2}}
\chapters@set{20261125T001000}{1}% zeros-to-ramp-sums
\chapters@set{20261125T001100}{2}% inert-primes-and-estimate-b
\chapters@set{20261125T005000}{3}% galois-fourier-and-averaging
\chapters@set{20261125T005500}{4}% columns-from-a-sphere
\chapters@set{20261125T000000}{5}% slope-method
\chapters@set{20261126T030100}{6}% a-twist-for-each-prime
\chapters@set{20261125T003000}{7}% the-chain-and-the-square-root
\chapters@set{20261125T010000}{8}% rank-lemma-and-zero-estimate
\chapters@set{20261125T010500}{9}% poincare-series-and-near-orthogonality
\chapters@set{20261124T223000}{10}% elliptic-curves-for-the-proof
\chapters@set{20261125T009500}{11}% fock-model-and-theta-intertwiners
\chapters@set{20261124T223100}{12}% complex-multiplication-and-inert-reduction
\chapters@set{20261125T008000}{13}% one-cm-determinant
\chapters@set{20261124T214500}{14}% supersingular-theta-divisibility
\chapters@set{20261124T214600}{15}% the-isotropic-plane
\chapters@set{20261125T013700}{16}% pairings-on-a-formal-group-of-height-two
\chapters@set{20261125T013500}{17}% hecke-nondegeneracy-on-the-tree
\chapters@set{20261126T030200}{18}% the-bound-is-attained
\chapters@set{20261125T009000}{19}% kronecker-limit-and-chowla-selberg
\chapters@set{20261125T002000}{20}% definite-theta-lifts
\chapters@set{20261125T015500}{21}% siegel-weil-in-the-smallest-case
\chapters@set{20261125T002100}{22}% definite-theta-lifts-and-norms
\chapters@set{20261125T016500}{23}% hecke-operators-and-theta-lifts
\chapters@set{20261125T016700}{24}% divisibility-along-the-chain
\chapters@set{20261125T004000}{25}% denominators-and-sha
\chapters@set{20261125T004100}{26}% tate-twists-and-sha
\chapters@set{20261126T020430}{27}% the-minus-part-by-gauss-sums
\chapters@set{20261125T014100}{28}% congruence-modules
\chapters@set{20261125T014500}{29}% eisenstein-congruences-and-selmer-classes
\chapters@set{20261125T014700}{30}% pseudocharacters-and-the-length-bound
\chapters@set{20261125T011000}{31}% galois-representations-black-box
\chapters@set{20261126T030300}{32}% the-data
\chapters@set{20261126T020210}{33}% bessel-poincare-series-and-their-gram-operators
\chapters@set{20261126T030400}{34}% the-seesaw-identity
\chapters@set{20261126T030500}{35}% divisibility-of-the-composite-lift
\chapters@set{20261126T030600}{36}% theorem-a-short-forms
\chapters@set{20261126T030700}{37}% the-determinant-of-a-chain
\chapters@set{20261126T030800}{38}% the-cost-of-one-lift
\chapters@set{20261126T030900}{39}% denominators-of-projections-to-lifts
\chapters@set{20261126T031000}{40}% twists-with-few-congruences
\chapters@set{20261126T031100}{41}% theorem-c-and-the-main-inequality
\providecommand\chapnumber[1]{%
  \@ifundefined{chapters@#1}%
    {\PackageError{chapters}{Sheet #1 is not listed in the index}%
      {List it in 20261125T012000--grh-workbook__guide.org\MessageBreak
       and run scripts/make-chapters.py.}}%
    {\edef\thechap{\@nameuse{chapters@#1}}}}
\makeatother
\chapnumber{<sheet ID>}
\makeatletter
\providecommand\chapters@set[2]{\@namedef{chapters@#1}{#2}}
\chapters@set{20261125T001000}{1}% zeros-to-ramp-sums
\chapters@set{20261125T001100}{2}% inert-primes-and-estimate-b
\chapters@set{20261125T005000}{3}% galois-fourier-and-averaging
\chapters@set{20261125T005500}{4}% columns-from-a-sphere
\chapters@set{20261125T000000}{5}% slope-method
\chapters@set{20261126T030100}{6}% a-twist-for-each-prime
\chapters@set{20261125T003000}{7}% the-chain-and-the-square-root
\chapters@set{20261125T010000}{8}% rank-lemma-and-zero-estimate
\chapters@set{20261125T010500}{9}% poincare-series-and-near-orthogonality
\chapters@set{20261124T223000}{10}% elliptic-curves-for-the-proof
\chapters@set{20261125T009500}{11}% fock-model-and-theta-intertwiners
\chapters@set{20261124T223100}{12}% complex-multiplication-and-inert-reduction
\chapters@set{20261125T008000}{13}% one-cm-determinant
\chapters@set{20261124T214500}{14}% supersingular-theta-divisibility
\chapters@set{20261124T214600}{15}% the-isotropic-plane
\chapters@set{20261125T013700}{16}% pairings-on-a-formal-group-of-height-two
\chapters@set{20261125T013500}{17}% hecke-nondegeneracy-on-the-tree
\chapters@set{20261126T030200}{18}% the-bound-is-attained
\chapters@set{20261125T009000}{19}% kronecker-limit-and-chowla-selberg
\chapters@set{20261125T002000}{20}% definite-theta-lifts
\chapters@set{20261125T015500}{21}% siegel-weil-in-the-smallest-case
\chapters@set{20261125T002100}{22}% definite-theta-lifts-and-norms
\chapters@set{20261125T016500}{23}% hecke-operators-and-theta-lifts
\chapters@set{20261125T016700}{24}% divisibility-along-the-chain
\chapters@set{20261125T004000}{25}% denominators-and-sha
\chapters@set{20261125T004100}{26}% tate-twists-and-sha
\chapters@set{20261126T020430}{27}% the-minus-part-by-gauss-sums
\chapters@set{20261125T014100}{28}% congruence-modules
\chapters@set{20261125T014500}{29}% eisenstein-congruences-and-selmer-classes
\chapters@set{20261125T014700}{30}% pseudocharacters-and-the-length-bound
\chapters@set{20261125T011000}{31}% galois-representations-black-box
\chapters@set{20261126T030300}{32}% the-data
\chapters@set{20261126T020210}{33}% bessel-poincare-series-and-their-gram-operators
\chapters@set{20261126T030400}{34}% the-seesaw-identity
\chapters@set{20261126T030500}{35}% divisibility-of-the-composite-lift
\chapters@set{20261126T030600}{36}% theorem-a-short-forms
\chapters@set{20261126T030700}{37}% the-determinant-of-a-chain
\chapters@set{20261126T030800}{38}% the-cost-of-one-lift
\chapters@set{20261126T030900}{39}% denominators-of-projections-to-lifts
\chapters@set{20261126T031000}{40}% twists-with-few-congruences
\chapters@set{20261126T031100}{41}% theorem-c-and-the-main-inequality
\providecommand\chapnumber[1]{%
  \@ifundefined{chapters@#1}%
    {\PackageError{chapters}{Sheet #1 is not listed in the index}%
      {List it in 20261125T012000--grh-workbook__guide.org\MessageBreak
       and run scripts/make-chapters.py.}}%
    {\edef\thechap{\@nameuse{chapters@#1}}}}
\makeatother
\chapnumber{20261125T004100}
\renewcommand{\thesection}{\thechap.\arabic{section}}

\theoremstyle{definition}
\newtheorem{problem}{Problem}[section]
\newtheorem{definition}{Definition}[section]
\theoremstyle{plain}
\newtheorem{theorem}{Theorem}[section]
\theoremstyle{remark}
\newtheorem*{leanremark}{In the Lean formalization}

\newcommand{\Z}{\mathbb{Z}}
\newcommand{\Q}{\mathbb{Q}}
\newcommand{\R}{\mathbb{R}}
\newcommand{\F}{\mathbb{F}}
\newcommand{\Zl}{\Z_\ell}
\newcommand{\Ql}{\Q_\ell}
\newcommand{\Oc}{\mathcal{O}}
\newcommand{\Gal}{\operatorname{Gal}}
\newcommand{\Hom}{\operatorname{Hom}}
\newcommand{\Cl}{\operatorname{Cl}}
\newcommand{\len}{\operatorname{len}}
\newcommand{\num}{\operatorname{num}}
\newcommand{\Sha}{\mathrm{Sha}}
\newcommand{\Shao}{\Sha^{\mathrm{odd}}}
\newcommand{\tl}[1]{\texttt{#1}}

\setlength{\emergencystretch}{2em}
\title{Odd Tate twists}
\author{}
\date{November 26, 2026}

\begin{document}
\maketitle
\refstepcounter{chap}\label{chap}

\paragraph{Prerequisites.}
This chapter continues Chapter~\ref{ds-chap}, \emph{Denominators and the Stickelberger element}, whose notation \(A\), \(A^{(i)}\), \(A^\pm\), \(\omega\), \(\sigma_a\) (Section~\ref{ds-sec:stickelberger}, with \(p = \ell\)), Theorem~\ref{ds-thm:HR} and Problem~\ref{ds-pr:irregular} we use.  Continuous cohomology of profinite groups \cite[Ch.~I; Ch.~II, \S7]{NSW}, inflation--restriction, and the vanishing of \(H^q(G, M)\), \(q \ge 1\), for finite \(G\) of order invertible on \(M\) \cite[Ch.~VII--VIII]{SeL}.  Decomposition and inertia groups \cite[Ch.~I, \S9]{Ne}, and the Artin isomorphism between the ideal class group of a number field and the Galois group of its Hilbert class field \cite[Ch.~VI, \S6]{Ne}.  The duality theorems for Galois cohomology of number fields are stated in the text with references.  In Sections~\ref{sec:twists} and~\ref{sec:split}, \(\ell\) is an odd prime; \(\num(r)\) is the numerator of a rational number \(r\) in lowest terms.

\section{Tate twists over \texorpdfstring{\(\Q\)}{Q}}\label{sec:twists}

In the previous chapter Bernoulli values gave the orders of odd eigenspaces of a class group.  Which Galois cohomology groups account for the numerators and denominators of zeta values?  Let \(\Q_S\) be the largest extension of \(\Q\) unramified outside \(\ell\) and \(\infty\), \(G_{\Q,\ell} := \Gal(\Q_S/\Q)\), and \(\chi_\ell \colon \Gal(\overline{\Q}/\Q) \to \Zl^\times\) the cyclotomic character, which factors through \(G_{\Q,\ell}\).  For \(k \in \Z\), \(\Zl(k)\) is \(\Zl\) with \(\sigma\) acting by \(\chi_\ell(\sigma)^k\), and similarly for \(\Ql(k)\), \(\Ql/\Zl(k)\) and \(\Z/\ell(k)\).  Let \(W_e := \Ql/\Zl(-e)\), and let \(G_\ell := \Gal(\overline{\Q}_\ell/\Q_\ell)\), mapped to \(G_{\Q,\ell}\) by restriction along a fixed embedding \(\overline{\Q} \to \overline{\Q}_\ell\).

\paragraph{Definition.}
The \emph{Tate--Shafarevich group of the twist \(e\)} is
\[
  \Sha_\ell(\Q, e) := \ker\bigl( H^1(G_{\Q,\ell}, W_e) \to H^1(G_\ell, W_e) \bigr) ,
\]
and \(h_{\ell,e}(\Q) \in \{0, 1, 2, \dots, \infty\}\) is its length as a \(\Zl\)-module.  For a number field \(E\) one defines \(\Sha_\ell(E, e)\) and \(h_{\ell,e}(E)\) in the same way, with the local Galois groups at all primes of \(E\) above \(\ell\).

\begin{problem}\label{pr:H0}
If \((\ell - 1) \nmid k\), then \(H^0(G_{\Q,\ell}, \Ql/\Zl(k)) = 0\).  If \((\ell-1) \mid k\) and \(k \ne 0\), it is cyclic of order \(\ell^{1 + v_\ell(k)}\).  The same holds for \(G_\ell\) in place of \(G_{\Q,\ell}\).
\end{problem}

By von Staudt and Clausen \cite[Thm.~5.10]{Wa}, the denominator of \(B_k\), \(k \ge 2\) even, is the product of the primes \(p\) with \((p-1) \mid k\).

\begin{problem}\label{pr:den-zeta}
For even \(k \ge 2\), the power of \(\ell\) in the denominator of \(\zeta(1-k) = -B_k/k\) is \(\lvert H^0(G_{\Q,\ell}, \Ql/\Zl(k)) \rvert\).
\end{problem}

For even \(k \ge 2\), the \(\ell\)-part of the denominator of \(\zeta(1-k)\) is the order of \(H^0(G_{\Q,\ell},\Ql/\Zl(k))\).  The \(\ell\)-part of its numerator will be the order of \(\Sha_\ell(\Q,k-1)\).

\begin{problem}\label{pr:unram-zero}
Let \(p \ne \ell\), \(e \ne 0\), and \(G_p := \Gal(\overline{\Q}_p/\Q_p)\).  Every unramified class in \(H^1(G_p, W_e)\) is zero.  Hence \(\Sha_\ell(\Q, e)\) is the group of classes in \(H^1(\Gal(\overline{\Q}/\Q), W_e)\) that vanish at every prime.
\end{problem}

Let \(F := \Q(\zeta_\ell)\).

\begin{problem}\label{pr:mod-l}
Suppose that \((\ell - 1) \nmid e\).  Then \(H^1(G_{\Q,\ell}, \Z/\ell(-e)) \to H^1(G_{\Q,\ell}, W_e)[\ell]\) is an isomorphism, and so is its analogue for \(G_\ell\).  Hence \(\Sha_\ell(\Q, e)[\ell]\) is the group of classes in \(H^1(G_{\Q,\ell}, \Z/\ell(-e))\) that vanish on \(G_\ell\).
\end{problem}

\begin{problem}\label{pr:to-classgroup}
Suppose that \((\ell - 1) \nmid e\).  Restriction to \(\Gal(\Q_S/F)\), followed by the Artin map, identifies the group of Problem~\ref{pr:mod-l} with the group of homomorphisms \(\phi \colon A \to \F_\ell\) such that \(\phi(\sigma_a z) = \omega(a)^{-e} \phi(z)\), that is, with \(\Hom(A^{(-e)}, \F_\ell)\).
\end{problem}

Problems~\ref{pr:mod-l} and~\ref{pr:to-classgroup} prove the following.

\begin{theorem}\label{thm:sha-class}
Let \((\ell - 1) \nmid e\).  Then \(\Sha_\ell(\Q, e)[\ell] \cong \Hom(A^{(-e)}, \F_\ell)\).  In particular \(\Sha_\ell(\Q, e) = 0\) if and only if \(A^{(-e)} = 0\).
\end{theorem}

\begin{problem}\label{pr:odd-e}
Let \(e \ge 1\) be odd.  Then \(\Sha_\ell(\Q, e) \ne 0\) if and only if \(\ell \mid \num(B_{e+1}/(e+1))\).  For example, \(\Sha_{691}(\Q, 11) \ne 0\) and \(\Sha_{37}(\Q, 31) \ne 0\).
\end{problem}

We state four inputs as black boxes.  By Soulé \cite{So}, \(H^2(G_{\Q,\ell}, \Zl(k))\) is finite for \(k \ge 2\).  By Poitou--Tate duality \cite[Ch.~VIII, \S6]{NSW}, \(\Sha_\ell(\Q, e)\) is dual to the kernel of \(H^2(G_{\Q,\ell}, \Zl(e+1)) \to H^2(G_\ell, \Zl(e+1))\), and by local duality \cite[Ch.~VII, \S2]{NSW} \(H^2(G_\ell, \Zl(k))\) is dual to \(H^0(G_\ell, \Ql/\Zl(1-k))\).  By Tate's global Euler characteristic formula \cite[Ch.~VIII, \S7]{NSW}, for every finite \(G_{\Q,\ell}\)-module \(M\) of \(\ell\)-power order, \(\lvert H^0(M) \rvert \lvert H^2(M) \rvert / \lvert H^1(M) \rvert = \lvert M^{c = 1} \rvert / \lvert M \rvert\), where \(c\) is a complex conjugation.  Since \(\ell\) is odd, \(G_{\Q,\ell}\) has \(\ell\)-cohomological dimension at most \(2\) \cite[Ch.~VIII, \S3]{NSW}, so \(H^3(G_{\Q,\ell}, M) = 0\) for every \(\ell\)-primary \(M\).

\begin{problem}\label{pr:sha2}
Let \(e \ge 1\) be odd.  Then \(H^2(G_\ell, \Zl(e+1)) = 0\), and \(\Sha_\ell(\Q, e)\) is dual to \(H^2(G_{\Q,\ell}, \Zl(e+1))\).
\end{problem}

\begin{problem}\label{pr:rank}
Let \(k \ge 2\).  The \(\Zl\)-module \(H^1(G_{\Q,\ell}, \Zl(k))\) has rank \(1\) if \(k\) is odd and rank \(0\) if \(k\) is even, and its torsion subgroup is isomorphic to \(H^0(G_{\Q,\ell}, \Ql/\Zl(k))\).
\end{problem}

The main conjecture of Iwasawa theory, proved by Mazur and Wiles \cite{MW}, gives the order; see \cite[Ch.~15]{Wa}.

\begin{theorem}[Mazur--Wiles]\label{thm:MW}
Let \(e \ge 1\) be odd.  Then
\[
  h_{\ell,e}(\Q) = v_\ell\bigl(\num(B_{e+1}/(e+1))\bigr) .
\]
\end{theorem}

\begin{problem}\label{pr:BK}
Let \(k \ge 2\) be even.  Then
\[
  \frac{\lvert H^2(G_{\Q,\ell}, \Zl(k)) \rvert}{\lvert H^1(G_{\Q,\ell}, \Zl(k)) \rvert} = \lvert \zeta(1-k) \rvert_\ell^{-1} .
\]
For \(k = 12\) and \(\ell = 691\) both sides are \(691\).
\end{problem}

This is the simplest case of the Bloch--Kato conjecture \cite{BK}.

\section{Odd characters}\label{sec:split}

Over \(\Q\) the odd twists are controlled by Bernoulli numbers.  How does this look over an abelian field, and what does it say about sums over all \(\ell\)?

Let \(E \subseteq \Q(\zeta_\ell)\) and \(\Delta_E := \Gal(E/\Q)\).  Then \(\Delta_E\) acts on \(\Sha_\ell(E, e)\), by \((\delta z)(g) := \tilde\delta\, z(\tilde\delta^{-1} g \tilde\delta)\) for a lift \(\tilde\delta\) of \(\delta\), which is the direct sum of its \(\eta\)-parts over the characters \(\eta\) of \(\Delta_E\), with values in \(\mu_{\ell-1} \subset \Zl^\times\).  Let \(\Sha_\ell(\Q, e; \eta)\) be defined as \(\Sha_\ell(\Q, e)\), with \(W_e\) replaced by \(W_e \otimes \eta^{-1}\).  A character \(\eta\) is \emph{critical for \(e\)} if \(\eta(-1) = (-1)^{e+1}\), and \emph{plus for \(e\)} otherwise.  The \emph{minus part} of \(\Sha_\ell(E, e)\) is the sum of the \(\eta\)-parts over the critical characters; it is the part on which complex conjugation acts by \((-1)^{e+1}\).

\begin{problem}\label{pr:induced}
The \(\eta\)-part of \(\Sha_\ell(E, e)\) is isomorphic to \(\Sha_\ell(\Q, e; \eta)\).  In particular \(\Sha_\ell(\Q, e)\) is the part fixed by \(\Delta_E\).
\end{problem}

\begin{problem}\label{pr:eta-rank}
Let \(\eta = \omega^j\) with \((\ell - 1) \nmid (e + j)\).  Then \(\Sha_\ell(\Q, e; \eta)[\ell] \cong \Hom(A^{(-e-j)}, \F_\ell)\), and \(\eta\) is critical for \(e\) if and only if \(-e-j\) is odd.
\end{problem}

\begin{problem}\label{pr:parity}
For every \(e \ne 0\), the minus part of \(\Sha_\ell(\Q(\zeta_\ell), e)\) has the same \(\ell\)-rank as \(A^-\).
\end{problem}

Over \(\Q(\zeta_\ell)\) the critical characters account for the \(\ell\)-rank of \(A^-\).  Over \(\Q\), for each fixed positive odd \(e\), one Bernoulli numerator controls all odd primes.

\begin{problem}\label{pr:one-number}
For a nonzero integer \(Y\), \(\sum_\ell v_\ell(Y) \log \ell \le \log \lvert Y \rvert\).  Hence, for odd \(e \ge 1\),
\[
  \sum_{\ell \text{ odd}} h_{\ell,e}(\Q) \log \ell \le \log \Bigl\lvert \num\Bigl(\frac{B_{e+1}}{e+1}\Bigr) \Bigr\rvert .
\]
\end{problem}

\begin{problem}\label{pr:bernoulli-size}
For \(k \ge 1\), \(\log \lvert \num(B_{2k}/2k) \rvert \le 2k \log(2k) + 3\).
\end{problem}

For the data of Chapter~\ref{da-chap}, the rewritten Lean proof \cite{GL} uses a CM subfield of a cyclotomic field and the twists \(0 \le e < 5r\), where \(r\) is the target rank of that chapter.  The counterpart of Problem~\ref{pr:one-number} there is Theorem~\ref{thm:oddsha} below.  It gives one positive integer \(X\) of controlled size such that, for every prime \(\ell\) and every coefficient ring \(\Oc\) allowed in the theorem, \(v_\Oc(X)\) bounds the \(\Oc\)-length of the module of cocycles of Definition~\ref{def:shaodd} at every level \(n\).  The rewritten Lean proof does not prove the main conjecture.

\section{Odd cocycles}\label{sec:odd}

In this section \(\ell\) is any prime.  Fix an embedding \(\overline{\Q} \subset \mathbb{C}\), and let \(c \in \Gal(\overline{\Q}/\Q)\) be complex conjugation.  Let \(E \subset \overline{\Q}\) be a CM field, \(G_E := \Gal(\overline{\Q}/E)\), and \(\chi := \chi_\ell|_{G_E}\).  Then \(c G_E c = G_E\) and \(\chi(cgc) = \chi(g)\).  For a place \(\mathfrak{P}\) of \(\overline{\Q}\) let \(D_{\mathfrak{P}} \supseteq I_{\mathfrak{P}}\) be its decomposition and inertia groups in \(G_E\).  Let \(\Oc\) be the ring of integers of a finite extension of \(\Ql\), with valuation \(v_\Oc\) normalized by \(v_\Oc(\pi) = 1\) on a uniformizer \(\pi\) and \(v_\Oc(0) = \infty\); for \(n \ge 0\) let \(\Oc_n := \Oc/\ell^n\Oc\), and write \(\Oc_n(-e)\) for this quotient with \(g \in G_E\) acting by \(\chi(g)^{-e}\).  We write \(\len\) for length over \(\Oc\), with values in \(\{0, 1, \dots, \infty\}\).

\begin{definition}\label{def:shaodd}
For integers \(e, n \ge 0\), \(\Shao_\ell(E, \Oc; e, n)\) is the \(\Oc\)-module of maps \(z \colon G_E \to \Oc_n\) such that:
\begin{enumerate}
\item[(a)] \(z\) is locally constant and \(z(gh) = z(g) + \chi(g)^{-e} z(h)\) for all \(g, h \in G_E\);
\item[(b)] \(z = 0\) on \(I_{\mathfrak{P}}\) for every \(\mathfrak{P}\) above a prime other than \(\ell\);
\item[(c)] for every \(\mathfrak{P}\) above \(\ell\): \(z = 0\) on \(I_{\mathfrak{P}}\) if \(e = 0\), and \(z = 0\) on \(D_{\mathfrak{P}} \cap \ker \chi\) if \(e > 0\);
\item[(d)] if \(e > 0\), then \(z(cgc) = -z(g)\) for all \(g \in G_E\).
\end{enumerate}
\end{definition}

These are cocycles, not classes: nothing is divided out.  In (c) the kernel is that of the \(\ell\)-adic character, not of its reduction modulo \(\ell^n\).  For \(e > 0\), the rewritten Lean proof \cite{GL} describes this module as a model for the odd part of \(\Sha^1(E,\Oc_n(-e))\) and records that the comparison with \(\Sha^1\) is not proved there.  Problems~\ref{pr:cob-odd} and~\ref{pr:c-action} explain the shape of the definition.

\begin{problem}\label{pr:cob-odd}
Let \(\ell\) be odd and \(e > 0\).  For \(b \in \Oc_n\), the coboundary \(z_b(g) := (\chi(g)^{-e} - 1) b\) satisfies (a)--(c), and it satisfies (d) only if \(z_b = 0\).
\end{problem}

\begin{problem}\label{pr:c-action}
For \(z\) satisfying (a)--(c) put \((c \cdot z)(g) := (-1)^e z(cgc)\).  Then \(c \cdot z\) satisfies (a)--(c) and \(c \cdot (c \cdot z) = z\).  For \(e > 0\), condition (d) is equivalent to \(c \cdot z = (-1)^{e+1}z\).  For odd \(\ell\) this is the sign of the minus part in Section~\ref{sec:split}.
\end{problem}

\paragraph{The twist \(e = 0\).}
Then (a) says that \(z\) is a continuous homomorphism, (d) is void, and (b), (c) together say that \(z\) vanishes on every inertia group.

\begin{problem}\label{pr:e0-hom}
The Artin map identifies \(\Shao_\ell(E, \Oc; 0, n)\) with \(\Hom(\Cl(E), \Oc_n)\).
\end{problem}

\begin{problem}\label{pr:hom-length}
Let \(C \cong \bigoplus_i \Z/m_i\Z\) be a finite abelian group.  Then
\[
  \len \Hom(C, \Oc_n) = \sum_i \min\bigl(v_\Oc(m_i),\, n\, v_\Oc(\ell)\bigr) \le v_\Oc(\lvert C \rvert) ,
\]
with equality for all large \(n\).
\end{problem}

\begin{problem}\label{pr:e0-example}
For \(E = \Q(\sqrt{-23})\), \(\ell = 3\) and \(\Oc = \Z_3\), the length of \(\Shao_3(E, \Z_3; 0, n)\) is \(\min(1, n)\); for \(\ell \ne 3\) and every \(\Oc\) it is \(0\).
\end{problem}

Let \(h_E := \lvert \Cl(E) \rvert\).  Problems~\ref{pr:e0-hom} and~\ref{pr:hom-length} prove the following.

\begin{theorem}\label{thm:e0}
Let \(E \subset \overline{\Q}\) be a CM field, \(\ell\) a prime, \(\Oc\) the ring of integers of a finite extension of \(\Ql\), and \(n \ge 0\).  Then \(\len \Shao_\ell(E, \Oc; 0, n) \le v_\Oc(h_E)\).
\end{theorem}

\paragraph{The field \(\Q(\zeta_\ell)\) modulo \(\ell\).}
Now let \(\ell\) be odd, \(E = F := \Q(\zeta_\ell)\), \(\Oc = \Zl\), \(n = 1\) and \(e > 0\), and let \(\lambda\) be the prime of \(F\) above \(\ell\).

\begin{problem}\label{pr:F-hom}
Every \(z \in \Shao_\ell(F, \Zl; e, 1)\) is a homomorphism \(G_F \to \F_\ell\).  The module \(\Shao_\ell(F, \Zl; e, 1)\) does not depend on \(e > 0\).
\end{problem}

\begin{problem}\label{pr:F-cyc}
The map \(z_{\mathrm{cyc}}(g) := (\chi(g) - 1)/\ell \bmod \ell\) is a homomorphism \(G_F \to \F_\ell\) satisfying (b) and (c), it is nonzero on \(D_{\mathfrak{P}}\) for every \(\mathfrak{P}\) above \(\ell\), and \(z_{\mathrm{cyc}}(cgc) = z_{\mathrm{cyc}}(g)\).  So \(z_{\mathrm{cyc}} \notin \Shao_\ell(F, \Zl; e, 1)\).
\end{problem}

\begin{problem}\label{pr:F-split}
Every \(z \in \Shao_\ell(F, \Zl; e, 1)\) vanishes on all of \(D_{\mathfrak{P}}\), for every \(\mathfrak{P}\) above \(\ell\).
\end{problem}

\begin{problem}\label{pr:F-minus}
The Artin map identifies \(\Shao_\ell(F, \Zl; e, 1)\) with \(\Hom(A^-, \F_\ell)\), for every \(e > 0\).  Hence its length is the \(\ell\)-rank of \(A^-\).
\end{problem}

By Problem~\ref{pr:parity}, this is also the \(\ell\)-rank of the minus part of \(\Sha_\ell(F, e)\).  So in this case the odd cocycles see what the minus part sees modulo \(\ell\): condition (d) removes the plus part, and (c) and (d) together remove the cyclotomic homomorphism \(z_{\mathrm{cyc}}\).

\paragraph{One integer for all \(\ell\).}
For a CM field in \(\Q(\zeta_N)\), the following theorem gives one integer for every prime \(\ell\), every coefficient ring specified below, every \(0 \le e < t\), and every \(n \ge 0\).

\begin{theorem}\label{thm:oddsha}
Let \(N \ge 1\) and \(t \ge 0\) be integers, let \(d_N\) be the exponent of the group \((\Z/N\Z)^\times\), and let \(E \subset \overline{\Q}\) be a CM field contained in \(\Q(\zeta_N)\).  There is an integer \(X\) with
\[
  1 \le X \le 2^{4(N + t + 3)^{27}} h_E
\]
such that for every prime \(\ell\) (including \(\ell = 2\)), every ring of integers \(\Oc\) of a finite extension of \(\Ql\) that contains a primitive \(d_N\)-th root of unity, every integer \(e\) with \(0 \le e < t\), and every \(n \ge 0\),
\[
  \len \Shao_\ell(E, \Oc; e, n) \le v_\Oc(X) .
\]
\end{theorem}

In particular these lengths are finite.  For \(e = 0\) the proof first applies Theorem~\ref{thm:e0} with right-hand side \(v_\Oc(h_E)\); when \(t = 1\) its chosen integer is \(X = h_E\).  For \(0 < e < t\) the rewriting first constructs an integer \(1 \le X_e \le 2^{(N+e+3)^6}\) satisfying a divisibility condition at every prime and a congruence for twisted Stickelberger sums.  Its proof of the bound then uses Jacobi-type quotients of Gauss sums and Kolyvagin's derivative cocycles, and treats \(\ell \le (N+t+3)^2\) and \(\ell > (N+t+3)^2\) separately.  Chapter~\ref{mi-chap} illustrates the Gauss-sum method.  The rewritten Lean proof describes the result as a finite-level upper bound of explicit size, and proves no equality of characteristic ideals.  For the field \(E\) of the data of Chapter~\ref{da-chap} it chooses a witness with \(t = 5r\), and the first term of its denominator error at \(\ell\) is \(t\, v_\ell(X)\).

\begin{leanremark}
\raggedright In the rewritten Lean proof \cite{GL}, Theorem~\ref{thm:oddsha} is \tl{New.exists\_length\_shaOdd\_le} and Theorem~\ref{thm:e0} is \tl{New.length\_shaOdd\_zero\_le\_addVal\_classNumber}, both in \tl{grh/src/Lib/Cyclotomic/OddSha.lean}, Definition~\ref{def:shaodd} is the case \(\Lambda = \Oc\) of \tl{New.Selmer.shaOdd} in \tl{grh/src/Lib/GaloisRep/Main/SelmerModules.lean}, and the application to the data uses \tl{New.Selmer.oddClassBound} and \tl{New.Selmer.denominatorError} in \tl{grh/src/New/Congruences/FirstOrder/TraceRelations.lean} and \tl{New.Selmer.isNegligible\_padicValNat\_oddClassBound} in \tl{grh/src/New/Congruences/FitTwistsExist.lean}.  The proof chooses \(X_e\) using \tl{New.OddPart.exists\_isStickelbergerBound}.  The proof of \tl{New.OddPart.exists\_isJacobiFamily} constructs the Jacobi-type quotients of Gauss sums, and the proof of \tl{New.OddPart.exists\_isKolyvaginSystem} constructs the derivative cocycles.  The introduction to \tl{OddSha.lean} gives the status of the comparison with \(\Sha^1\) and of the finite-level bound.
\end{leanremark}

\begin{problem}\label{pr:37}
Let \(E = \Q(\zeta_{37})\), \(N = 37\) and \(t \ge 2\).  Then \(\Z_{37}\) satisfies the hypotheses of Theorem~\ref{thm:oddsha}, and every integer \(X\) as in that theorem is divisible by \(37\).
\end{problem}

\begin{problem}\label{pr:sum-ell}
Let \(N\), \(t\), \(d_N\), \(E\), \(X\) be as in Theorem~\ref{thm:oddsha}, and for each prime \(\ell\) let \(\Oc_{(\ell)}\) be the ring of integers of \(\Ql(\zeta_{d_N})\).  Then \(\Oc_{(\ell)}\) satisfies the hypotheses of the theorem, and for \(0 \le e < t\), \(n \ge 0\) and every finite set \(T\) of primes,
\[
  \sum_{\ell \in T} \frac{\len \Shao_\ell(E, \Oc_{(\ell)}; e, n)}{v_{\Oc_{(\ell)}}(\ell)} \log \ell \le \log X \le 4 (N + t + 3)^{27} \log 2 + \log h_E .
\]
\end{problem}

This is the analogue of Problem~\ref{pr:one-number}.  In the application to Chapter~\ref{da-chap}, the rewritten Lean proof sums \(v_\ell(X)\log\ell\) over a finite set of primes and uses the size bound for \(X\) in Theorem~\ref{thm:oddsha}.

\section*{Hints}

\subsection*{\S\ref{sec:twists}}
\begin{description}
\item[\ref{pr:H0}] \(\chi_\ell\) maps \(G_{\Q,\ell}\) and \(G_\ell\) onto \(\Zl^\times\).  Compute the least \(v_\ell(u^k - 1)\) over \(u \in \Zl^\times\).
\item[\ref{pr:den-zeta}] Combine Problem~\ref{pr:H0} with the theorem of von Staudt and Clausen.  If \((\ell - 1) \nmid k\), then \(B_k/k\) is \(\ell\)-integral \cite[Cor.~5.14]{Wa}.
\item[\ref{pr:unram-zero}] An unramified class lies in \(H^1(\hat{\Z}, W_e) = W_e/(\varphi - 1)W_e\), where Frobenius \(\varphi\) acts by \(p^{-e}\).  For the second statement, inflate from \(G_{\Q,\ell}\); conversely, by inflation--restriction a class that is unramified outside \(\ell\) comes from \(G_{\Q,\ell}\), since \(\Gal(\overline{\Q}/\Q_S)\) is generated as a closed normal subgroup by the inertia groups at primes \(p \ne \ell\).
\item[\ref{pr:mod-l}] Use the sequence \(0 \to \Z/\ell(-e) \to W_e \xrightarrow{\ell} W_e \to 0\) and Problem~\ref{pr:H0}.
\item[\ref{pr:to-classgroup}] Use inflation--restriction to \(G_F := \Gal(\Q_S/F)\), which acts trivially on \(\Z/\ell(-e)\).  The prime above \(\ell\) in \(F\) is principal, and \([F_\lambda : \Q_\ell] = \ell - 1\) is prime to \(\ell\).
\item[\ref{pr:odd-e}] Let \(i \equiv -e \pmod{\ell - 1}\).  Use Theorem~\ref{thm:sha-class}, Theorem~\ref{ds-thm:HR}, and Kummer's congruence \(B_m/m \equiv B_n/n \pmod \ell\) for even \(m, n \ge 2\) with \(m \equiv n \not\equiv 0 \pmod{\ell-1}\) \cite[Ch.~5]{Wa}; treat \(i \equiv 1\) separately, using \(A^{(1)} = 0\) \cite[Prop.~6.16]{Wa}.
\item[\ref{pr:sha2}] Problem~\ref{pr:H0} with \(k = -e\), and local duality.
\item[\ref{pr:rank}] Apply the Euler characteristic formula to \(\Zl(k)/\ell^n\) and let \(n \to \infty\); since the groups \(H^i(G_{\Q,\ell}, \Zl(k)/\ell^n)\) are finite, their inverse limit is \(H^i(G_{\Q,\ell}, \Zl(k))\) \cite[Ch.~II, \S7]{NSW}.  For the upper bound on the rank, show that \(H^2(G_{\Q,\ell}, \Zl(k)/\ell^n) = H^2(G_{\Q,\ell}, \Zl(k))/\ell^n\), which is bounded by Soulé's theorem.  For the torsion, use \(0 \to \Zl(k) \to \Ql(k) \to \Ql/\Zl(k) \to 0\).  The groups \(H^i(G_{\Q,\ell}, M)\) are finite for finite \(G_{\Q,\ell}\)-modules \(M\) of \(\ell\)-power order \cite[Ch.~VIII, \S3]{NSW}; finite generation of \(H^1(G_{\Q,\ell}, \Zl(k))\) follows from finiteness of its reduction modulo \(\ell\) by Nakayama's lemma, and \(H^3 = 0\) passes to the limit.
\item[\ref{pr:BK}] Problems~\ref{pr:sha2}, \ref{pr:rank} and~\ref{pr:den-zeta}, and Theorem~\ref{thm:MW} with \(e = k - 1\).
\end{description}

\subsection*{\S\ref{sec:split}}
\begin{description}
\item[\ref{pr:induced}] Shapiro's lemma for \(\mathrm{Ind}_{G_{E,\ell}}^{G_{\Q,\ell}} W_e \cong \bigoplus_\eta W_e \otimes \eta\).  The field \(E\) has one prime above \(\ell\).
\item[\ref{pr:eta-rank}] Repeat Problems~\ref{pr:mod-l} and~\ref{pr:to-classgroup} with the twisted coefficients.  On \((W_e\otimes\eta^{-1})[\ell]\), the action is through the reduction of \(\omega^{-e-j}\).
\item[\ref{pr:parity}] Sum Problem~\ref{pr:eta-rank} over the critical \(j\).  The one character \(\omega^{-e}\) not covered there is plus for \(e\).
\item[\ref{pr:one-number}] For the second part, Theorem~\ref{thm:MW}.
\item[\ref{pr:bernoulli-size}] \(B_{2k} = (-1)^{k+1} 2 (2k)! \zeta(2k)/(2\pi)^{2k}\), and bound the denominator by von Staudt--Clausen and \(\prod_{p \le n} p \le 4^n\).
\end{description}

\subsection*{\S\ref{sec:odd}}
\begin{description}
\item[\ref{pr:cob-odd}] \(\chi\) is trivial on inertia at primes other than \(\ell\), and \(\chi(cgc) = \chi(g)\); \(2\) is a unit in \(\Oc_n\).
\item[\ref{pr:c-action}] \(\chi_\ell(c) = -1\).  Conjugation by \(c\) carries \(D_{\mathfrak{P}}\), \(I_{\mathfrak{P}}\) to \(D_{c\mathfrak{P}}\), \(I_{c\mathfrak{P}}\).
\item[\ref{pr:e0-hom}] A homomorphism killing all inertia groups factors through the Galois group of the Hilbert class field; \(E\) has no real places.
\item[\ref{pr:hom-length}] \(\Hom(\Z/m\Z, \Oc_n) \cong \Oc/(m, \ell^n)\), and \(\Oc/\pi^k\Oc\) has length \(k\).
\item[\ref{pr:e0-example}] The class number of \(\Q(\sqrt{-23})\) is \(3\).
\item[\ref{pr:F-hom}] \(\chi(g) \equiv 1 \pmod \ell\) on \(G_F\).
\item[\ref{pr:F-cyc}] \(\chi(D_{\mathfrak{P}}) = 1 + \ell\Zl\), since \(\Q_\ell(\zeta_{\ell^2})\) has degree \(\ell\) over \(\Q_\ell(\zeta_\ell)\).
\item[\ref{pr:F-split}] On \(D_{\mathfrak{P}}\), \(z\) factors through \(\chi(D_{\mathfrak{P}}) = 1 + \ell\Zl\), so it is a multiple of \(z_{\mathrm{cyc}}\) there.  Compare at \(\mathfrak{P}\) and \(c\mathfrak{P}\), which are conjugate under \(G_F\).
\item[\ref{pr:F-minus}] Problems~\ref{pr:F-split} and~\ref{pr:c-action}; \(\lambda = (1 - \zeta_\ell)\) is principal, the Artin map satisfies \(\mathrm{Art}(\sigma \mathfrak{a}) = \tilde\sigma\, \mathrm{Art}(\mathfrak{a})\, \tilde\sigma^{-1}\), and \(c = \sigma_{-1}\) acts on \(A^{(i)}\) by \((-1)^i\).
\item[\ref{pr:37}] Problems~\ref{pr:F-minus} and~\ref{ds-pr:irregular}.
\item[\ref{pr:sum-ell}] \(v_\Oc(X) = v_\Oc(\ell)\, v_\ell(X)\), and Problem~\ref{pr:one-number}.
\end{description}

\begin{thebibliography}{NSW}
\bibitem[BK]{BK} S.~Bloch and K.~Kato, \emph{\(L\)-functions and Tamagawa numbers of motives}, in: The Grothendieck Festschrift, Vol.~I, Progr.\ Math.\ 86, Birkh\"auser, 1990, 333--400.
\bibitem[GL]{GL} L.~Sernau, \emph{GRH for Dirichlet \(L\)-functions, in Lean~4}, \url{https://github.com/lsernau/grh-lean}.
\bibitem[MW]{MW} B.~Mazur and A.~Wiles, \emph{Class fields of abelian extensions of \(\mathbf{Q}\)}, Invent.\ Math.\ \textbf{76} (1984), 179--330.
\bibitem[Ne]{Ne} J.~Neukirch, \emph{Algebraic Number Theory}, Grundlehren Math.\ Wiss.\ 322, Springer, 1999.
\bibitem[NSW]{NSW} J.~Neukirch, A.~Schmidt and K.~Wingberg, \emph{Cohomology of Number Fields}, 2nd ed., Grundlehren Math.\ Wiss.\ 323, Springer, 2008.
\bibitem[SeL]{SeL} J.-P.~Serre, \emph{Local Fields}, Grad.\ Texts in Math.\ 67, Springer, 1979.
\bibitem[So]{So} C.~Soul\'e, \emph{\(K\)-th\'eorie des anneaux d'entiers de corps de nombres et cohomologie \'etale}, Invent.\ Math.\ \textbf{55} (1979), 251--295.
\bibitem[Wa]{Wa} L.~C.~Washington, \emph{Introduction to Cyclotomic Fields}, 2nd ed., Grad.\ Texts in Math.\ 83, Springer, 1997.
\end{thebibliography}

\end{document}
